\section{Two Sets II}
\label{sec:cses-two-sets-ii}

\problemstatement{Count the number of ways to partition the numbers
$\{1,2,\ldots,n\}$ into two subsets with \emph{equal} sums. Two partitions
that are reflections of each other count as the same. Output modulo
$10^9+7$.}

\begin{patternbox}
\emph{``Split $\{1..n\}$ into two equal-sum halves, count the ways''} is a
\textbf{counting subset-sum}. The total is
$T=\frac{n(n+1)}{2}$; if $T$ is odd the answer is $0$, else we count
subsets summing to $T/2$, then divide by $2$ to remove the mirror
duplication.
\end{patternbox}

\subsection*{State and transition}

Let the target be $t=T/2$. Define $\textit{dp}[s]$ = number of subsets of
$\{1,\ldots,n\}$ (processed one number at a time) that sum to $s$, with
$\textit{dp}[0]=1$. Each number $k$ is either in the subset or not -- a
$0/1$ choice -- so with the downward inner loop:
\[
  \textit{dp}[s]\mathrel{+}=\textit{dp}[s-k]\quad\text{for } s=t\ \text{down to}\ k.
\]
Each valid partition is counted twice (subset $A$ and its complement
$B$), so the final answer is $\textit{dp}[t]/2$ modulo $10^9+7$ (using the
modular inverse of $2$).

\begin{ideabox}[Count Subsets, Then Halve for the Mirror]
Counting subsets that reach exactly half the total counts every partition
from both sides. Dividing by $2$ (via the modular inverse, since we are
mod a prime) removes that symmetry. The odd-total shortcut avoids
pointless work when no equal split exists.
\end{ideabox}

\subsection*{Algorithm}

\begin{enumerate}
  \item Compute $T=n(n+1)/2$; if $T$ is odd, output $0$.
  \item Set $t=T/2$, $\textit{dp}[0]=1$. For $k=1\ldots n-1$ (one number
        may be fixed to a side to pre-halve, or process all and divide):
        update $\textit{dp}[s]\mathrel{+}=\textit{dp}[s-k]$ downward.
  \item Output $\textit{dp}[t]$, halved via the modular inverse of $2$.
\end{enumerate}

(Alternatively, process $k=1\ldots n-1$ only -- fixing $n$ to one side --
which directly avoids the double count without a division.)

\subsection*{Dry run}

\dryrun{$n=7$, $T=28$, target $t=14$.}

\begin{itemize}
  \item Subsets of $\{1..7\}$ summing to $14$ include $\{7,6,1\},
        \{7,5,2\},\{7,4,3\},\{7,4,2,1\},\ldots$
  \item The counting DP gives $\textit{dp}[14]=8$; halving yields the
        answer $4$ distinct partitions.
\end{itemize}

\subsection*{C++ implementation}

\begin{lstlisting}
#include <bits/stdc++.h>
using namespace std;
const long long MOD = 1e9 + 7;

int main() {
    int n;
    cin >> n;
    long long total = (long long)n * (n + 1) / 2;
    if (total % 2 != 0) { cout << 0 << "\n"; return 0; }

    int t = total / 2;
    vector<long long> dp(t + 1, 0);
    dp[0] = 1;
    // fix element n on one side: process 1..n-1, no division needed
    for (int k = 1; k <= n - 1; k++)
        for (int s = t; s >= k; s--)
            dp[s] = (dp[s] + dp[s - k]) % MOD;

    cout << dp[t] << "\n";
    return 0;
}
\end{lstlisting}

\begin{complexitybox}
\textbf{Time Complexity.} $O(n\cdot T)=O(n^3)$ in the worst sense, but
$T=O(n^2)$ so it is $O(n\cdot n^2)$; practically $O(n\cdot t)$ with
$t\approx n^2/4$. \textbf{Space Complexity.} $O(t)=O(n^2)$.
\end{complexitybox}

\begin{takeawaybox}
Two Sets II is counting subset-sum with a symmetry fix. The clean way to
avoid the double count is to fix one fixed element (here $n$) to a single
side and only distribute the rest -- turning ``divide by two'' into ``count
directly.'' Recognise equal-partition problems as subset-sum to half the
total.
\end{takeawaybox}
